% !TEX root = _main.tex
% ========================================
% 作品タイトル：「日暮どきこそ耳を澄まして：「Echo」〜音だけの動物物語〜」
% =========================================
\clearpage
\onecolumn

\subsection*{付録A：音響イベント一覧（Event 1–149 全網羅表）}
\addcontentsline{toc}{subsection}{付録A：音響イベント一覧（Event 1–149 全網羅表）}
\label{app:all-events}


% 行テンプレ（必要なら列を増減）
% 例：\EventRow{1}{第1場面}{環境音}{-}{雷雨}{持続}{序盤の切迫感と一貫性}
\newcommand{\EventRow}[7]{#1 & #2 & #3 & #4 & #5 & #6 & #7\\}

\setlength{\LTpre}{0pt}
\setlength{\LTpost}{0pt}
\renewcommand{\arraystretch}{1.15}

\scriptsize

\begin{longtable}{@{}r p{1.0cm} p{1.0cm} p{2.0cm} p{2.8cm} p{3.0cm} p{6.1cm}@{}}
	\caption{音響イベント全一覧（Event 1–149）}\label{tab:all-events}                                                  \\
	\toprule
	\textbf{ID} & \textbf{場面} & \textbf{分類} & \textbf{主体} & \textbf{素材名} & \textbf{操作} & \textbf{役割・意図（要約）} \\
	\midrule
	\endfirsthead

	\toprule
	\textbf{ID} & \textbf{場面} & \textbf{分類} & \textbf{主体} & \textbf{素材名} & \textbf{操作} & \textbf{役割・意図（要約）} \\
	\midrule
	\endhead

	\midrule
	\multicolumn{7}{r}{\footnotesize 次ページに続く}                                                               \\
	\endfoot

	\bottomrule
	\endlastfoot

	% --------------------
	% ここに Event 行を並べる（1〜149）
	% --------------------
	\EventRow{1}{第1場面}{環境音}{環境}{雷雨}{持続（背景）}{墜落直前の嵐を提示し，切迫感と場面の一貫性を支える}
	\EventRow{2}{第1場面}{鳴き声}{ハチドリ（群れ）}{ハチドリ鳴き声（素材A）}{重ね+リバーブ付加}{俯瞰視点で群れの混沌と一体感・距離感を表現}
	\EventRow{3}{第1場面}{鳴き声}{ハチドリ（群れ）}{ハチドリ鳴き声（素材B）}{重ね+リバーブ付加}{俯瞰視点で群れの混沌と一体感・距離感を表現}
	\EventRow{4}{第1場面}{鳴き声}{ハチドリ（群れ）}{ハチドリ鳴き声（素材C）}{重ね+リバーブ付加}{俯瞰視点で群れの混沌と一体感・距離感を表現}
	\EventRow{5}{第1場面}{鳴き声}{ハチドリ（群れ）}{ハチドリ鳴き声（素材D）}{重ね+リバーブ付加}{俯瞰視点で群れの混沌と一体感・距離感を表現}
	\EventRow{6}{第1場面}{鳴き声}{ハチドリ（群れ）}{ハチドリ鳴き声（素材E）}{重ね+リバーブ付加}{俯瞰視点で群れの混沌と一体感・距離感を表現}
	\EventRow{7}{第1場面}{環境音}{環境（主人公近傍）}{風切り（明瞭な風）}{追加レイヤ／近接感強調}{主人公視点への切替を示し，身体的近接感を提示}
	\EventRow{8}{第1場面}{鳴き声}{ハチドリ（主）}{ハチドリ鳴き声（上昇要素）}{ドライ（リバーブなし）}{主人公が仲間を呼ぶ／不安・疑問のニュアンスを提示}
	\EventRow{9}{第1場面}{鳴き声}{仲間ハチドリ}{ハチドリ鳴き声（落ち着き）}{リバーブ付加}{仲間が一定距離にいることを示し，頼りとなる存在を提示}
	\EventRow{10}{第1場面}{鳴き声}{ハチドリ（主）}{ハチドリ鳴き声（切迫）}{強め素材選択}{緊急性・高揚を示し，次の急行動への予兆}
	\EventRow{11}{第1場面}{効果音}{ハチドリ（主）}{突風／体勢転換の風切り}{瞬間的・勢い強調}{全力飛行への移行点を示し，運動の急変を強調}
	\EventRow{12}{第1場面}{効果音}{ハチドリ（主）}{急加速の風切り（別素材）}{重ね配置+音量−28$\to$−6 dBランプ}{速度上昇と衝突の迫りを音量変化で表現し緊迫感を増幅}
	\EventRow{13}{第1場面}{鳴き声}{ハチドリ（主）}{ハチドリ鳴き声（更に切迫）}{素材選択（切迫強）}{焦燥・危機感の増大を示し，緊張を上積み}
	\EventRow{14}{第1場面}{効果音}{ハチドリ（主）}{衝突音}{リバーブ付加}{視覚なしで衝突を明確化し，衝撃の大きさと余韻を強調}
	\EventRow{15}{第1場面}{鳴き声}{ハチドリ（主）}{墜落時の鳴き声}{同一素材反復+ベロシティ128$\to$52$\to$45$\to$28$\to$15$\to$6$\to$1}	{疑似ディレイで墜落の時間的引き伸ばしと注意喚起を行う}
	\EventRow{16}{第1場面}{効果音}{ハチドリ（主）}{着地衝突音}{音量−1.6 dB（抑制）}{致命的でない着地／衝撃緩和を示す}
	\EventRow{17}{第2場面}{環境音}{環境}{雨（雷なし）}{素材選択（雫明瞭）}{地上にいる状況を示しつつ，不安が残る環境を提示}
	\EventRow{18}{第2場面}{足音}{ウサギ}{芝生の上を歩く}{タイミング同期+ベロシティ上昇（接近）}{跳ね移動の着地（硬め）成分を表現し，接近を提示}
	\EventRow{19}{第2場面}{足音}{ウサギ}{草むらを歩く}{タイミング同期+ベロシティ上昇（接近）}{跳ね移動の草揺れ成分を表現し，接近を提示}
	\EventRow{20}{第2場面}{鳴き声}{ウサギ}{ウサギ鳴き声（短間隔反復）}{反復配置}{敵意でなく様子見の存在であることを直感的に示す}
	\EventRow{21}{第2場面}{鳴き声}{ハチドリ（主）}{ハチドリ鳴き声（短音2回）}{ノート極短+ピッチ上行（驚き）／下行（落胆）}{起き上がり$\to$倒れ込みの反応と心理（驚き／疲弊）を示す}
	\EventRow{22}{第2場面}{効果音}{ハチドリ（主）}{草むらを歩く}{鳴き声と同時再生}{鳴き声だけでは不足する身体動作（起きる／倒れる）を補完}
	\EventRow{23}{第2場面}{鳴き声}{ウサギ}{ウサギ鳴き声}{予兆として配置}{主人公を運ぶ行動に入る前触れを示す}
	\EventRow{24}{第2場面}{鳴き声}{ウサギ}{ウサギ鳴き声}{G2$\to$G3（溜め$\to$掛け声）}{持ち上げ動作の力の込め方を聴覚的に示す}
	\EventRow{25}{第2場面}{鳴き声}{ハチドリ（主）}{ハチドリ鳴き声（短音2回）}{極短ノート+C3$\to$A3（上行）}{ウサギの行動への驚きを疑問調の上昇で提示}
	\EventRow{26}{第2場面}{効果音}{ウサギ}{手で穴を掘る1}{二段階音を動作に対応}{主人公を乗せて持ち上げる身体動作を直感的に強調}
	\EventRow{27}{第2場面}{足音}{ウサギ}{芝生の上を歩く}{タイミング同期+ベロシティ下降（遠ざかり）}{去っていく跳ね移動の着地成分を表現し距離感を提示}
	\EventRow{28}{第2場面}{足音}{ウサギ}{草むらを歩く}{タイミング同期+ベロシティ下降（遠ざかり）}{去っていく跳ね移動の草揺れ成分を表現し距離感を提示}
	\EventRow{29}{第3場面}{環境音}{環境（巣穴外）}{街の小雨}{強めリバーブ付加}{雨が遮蔽された空間から間接的に聴こえる距離・空間を表現}
	\EventRow{30}{第3場面}{環境音}{環境（巣穴内）}{下水道（滴り）}{素材選択}{巣穴内部の水滴環境を示し，屋外雨と質感を分ける}
	\EventRow{31}{第3場面}{鳴き声}{ハチドリ（主）}{ハチドリ鳴き声（短音2回）}{極短ノート+A2$\to$B3}{目覚めの戸惑いと状況把握できない反応を示す}
	\EventRow{32}{第3場面}{効果音}{ウサギ}{rabbit-eating}{同時配置}{落ち着いたウサギを提示し，主人公の戸惑いと対照を作る}
	\EventRow{33}{第3場面}{効果音}{ハチドリ（主）}{草むらを歩く}{同時配置}{起き上がり動作を補完し，身体状態を明示}
	\EventRow{34}{第3場面}{鳴き声}{ウサギ}{ウサギ鳴き声}{区切りとして配置}{食事終了と次行動への予兆を示す}
	\EventRow{35}{第3場面}{鳴き声}{ウサギ}{ウサギ鳴き声}{B2$\to$F3（溜め$\to$掛け声）}{食べ物を転がす動作の意図と力点を示す}
	\EventRow{36}{第3場面}{効果音}{ウサギ（行為）}{ビリヤードでポケットする1}{素材転用}{硬質な食べ物が転がる音を代替し，硬さも示す}
	\EventRow{37}{第3場面}{鳴き声}{ハチドリ（主）}{ハチドリ鳴き声}{F2$\to$C$\sharp$2$\to$A2（下げて上げる）}{思考の揺らぎを含む戸惑いを表現}
	\EventRow{38}{第3場面}{鳴き声}{ウサギ}{ウサギ鳴き声}{短く単一音}{落ち着いた返答で主人公と対照を作る}
	\EventRow{39}{第3場面}{鳴き声}{ハチドリ（主）}{ハチドリ鳴き声}{B2$\to$F2（下降）}{暫定的理解+戸惑い残存を下降で表現}
	\EventRow{40}{第3場面}{鳴き声}{ウサギ}{ウサギ鳴き声}{D3（高め）}{肯定的態度で「食べてよい」を許可する役割}
	\EventRow{41}{第3場面}{鳴き声}{ハチドリ（主）}{ハチドリ鳴き声}{短い単音}{次の咀嚼音の主体が主人公であることを予告}
	\EventRow{42}{第3場面}{効果音}{ハチドリ（主）}{ガムを噛む}{不規則配置}{自然な咀嚼のリズムを作り，摂食を明示}
	\EventRow{43}{第3場面}{効果音}{ハチドリ（主）}{飲む}{直後配置}{飲み込み動作の完了を示す}
	\EventRow{44}{第3場面}{鳴き声}{ハチドリ（主）}{ハチドリ鳴き声}{F3$\to$E3$\to$D2$\to$C3}{食後の疲弊を残しつつ動作完了を示し，後続への連続性を確保}
	\EventRow{45}{第3場面}{鳴き声}{ウサギ}{ウサギ鳴き声}{F3（高め）}{食事成功への喜び／肯定を示す}
	\EventRow{46}{第3場面}{鳴き声}{ハチドリ（主）}{ハチドリ鳴き声}{C3$\to$A3（瞬間）}{突発刺激への反射的動揺を示し，第1場面を想起}
	\EventRow{47}{第3場面}{環境音}{環境}{落雷}{同時配置}{雷が動揺の原因であることを明確化}
	\EventRow{48}{第3場面}{鳴き声}{ハチドリ（主）}{ハチドリ鳴き声（短）}{短音}{動揺後の一時的収束（完全な落ち着きではない）}
	\EventRow{49}{第3場面}{鳴き声}{ウサギ}{ウサギ鳴き声}{F2（低め）}{心配・落ち着いた態度を低ピッチで表現}
	\EventRow{50}{第3場面}{鳴き声}{ハチドリ（主）}{ハチドリ鳴き声}{A2$\to$C2（瞬間）}{疲弊の再提示と落雷後の緊張余韻を形成}
	\EventRow{51}{第4場面}{環境音}{環境}{スズメの鳴き声}{環境扱い+リバーブ付加}{翌朝の平穏を提示し，雨との対比で転換を作る}
	\EventRow{52}{第4場面}{鳴き声}{ウサギ}{ウサギ鳴き声}{足音と同時}{移動主体がウサギであることを示す}
	\EventRow{53}{第4場面}{足音}{ウサギ}{芝生の上を歩く}{タイミング同期（一定）}{跳ね移動の着地成分を表現（距離変化なし）}
	\EventRow{54}{第4場面}{足音}{ウサギ}{草むらを歩く}{タイミング同期（一定）}{跳ね移動の草揺れ成分を表現（距離変化なし）}
	\EventRow{55}{第4場面}{鳴き声}{ハチドリ（主）}{ハチドリ鳴き声}{B2長め$\to$間$\to$G2}{落ち着いた声掛けで回復を示し，主人公の存在を明示}
	\EventRow{56}{第4場面}{鳴き声}{ウサギ}{ウサギ鳴き声}{D2}{主人公への肯定的返答を示す}
	\EventRow{57}{第4場面}{効果音}{環境}{風に揺れる草木2}{素材提示（強め）}{新たな気配（シカ出現）への前触れを作る}
	\EventRow{58}{第4場面}{鳴き声}{ハチドリ（主）}{ハチドリ鳴き声}{F2$\to$C$\sharp$2$\to$A2}{気配への戸惑い（立ち止まり$\to$反応）を表現}
	\EventRow{59}{第4場面}{鳴き声}{ウサギ}{ウサギ鳴き声}{F3（単調）}{注意を向けつつ落ち着いた態度の一貫性を保つ}
	\EventRow{60}{第4場面}{足音}{シカ（代替）}{馬が走る1}{C$\sharp$2+音量−20$\to$−2.1 dBランプ}{大型動物の接近を代替素材で表現し距離変化を提示}
	\EventRow{61}{第4場面}{鳴き声}{シカ}{baby-deer-calling-mama}{素材提示}{足音の主体がシカであることを確定させる}
	\EventRow{62}{第4場面}{鳴き声}{ウサギ}{ウサギ鳴き声}{A2$\to$F3（最初長め）}{突然の存在への困惑を示しつつ落ち着きも保持}
	\EventRow{63}{第4場面}{鳴き声}{シカ}{baby-deer-calling-mama（別箇所）}{素材箇所変更}{次の接近（64）の合図として主体を予告}
	\EventRow{64}{第4場面}{足音}{シカ（代替）}{馬が走る1}{C$\sharp$2+音量−2.1$\to$+6 dBランプ}{さらに近づく接近を音量で強調し速度低下も示す}
	\EventRow{65}{第4場面}{鳴き声}{ハチドリ（主）}{ハチドリ鳴き声（短）}{素材内ピッチ変化利用}{驚きはあるが過度ではない反応を単調回避しつつ提示}
	\EventRow{66}{第4場面}{鳴き声}{ウサギ}{ウサギ鳴き声}{E3（やや高め）}{落ち着きつつ状況把握できない様子を示す}
	\EventRow{67}{第4場面}{鳴き声}{ハチドリ（主）}{ハチドリ鳴き声（高$\to$低）}{素材選択}{嗅がれる驚きを過度な動揺なしで表現}
	\EventRow{68}{第4場面}{鳴き声}{ウサギ}{ウサギ鳴き声}{B2}{動揺は抑えつつ身構える反応を示す}
	\EventRow{69}{第4場面}{効果音}{シカ（代替）}{イノシシが鼻をフンフン}{C3$\to$間$\to$A2（後半低ピッチ）}{嗅ぎ取り動作を代替素材で表現し慎重な再確認も示す}
	\EventRow{70}{第4場面}{鳴き声}{シカ}{baby-deer-calling-mama（同ピッチ）}{ベロシティ117}{緊張を緩め，危険でない存在として印象を誘導}
	\EventRow{71}{第4場面}{鳴き声}{ハチドリ（主）}{ハチドリ鳴き声}{A3}{警戒がわずかに和らいだ反応を示す}
	\EventRow{72}{第4場面}{足音}{シカ（代替）}{馬が走る1}{C$\sharp$2+音量+6$\to$−2 dB}{道案内として先行し遠ざかる距離変化を表現}
	\EventRow{73}{第4場面}{鳴き声}{シカ}{baby-deer-calling-mama}{歩行後に間を空けて配置}{立ち止まり呼びかけ（ついて来い）の表現}
	\EventRow{74}{第4場面}{鳴き声}{ウサギ}{ウサギ鳴き声}{F2$\to$E3（極短）}{即時的戸惑いを短音と音高変化で表現}
	\EventRow{75}{第4場面}{鳴き声}{ハチドリ（主）}{ハチドリ鳴き声（短）}{A2$\to$F$\sharp$3（瞬間）}{ウサギに同調した戸惑い（強い動揺には至らない）}
	\EventRow{76}{第4場面}{環境音}{環境}{風に揺れる草木1}{持続配置}{沈黙区間を無音にせず場の空気感を保持}
	\EventRow{77}{第4場面}{鳴き声}{シカ}{baby-deer-calling-mama}{E3}{敵意のなさを高めのピッチで示す}
	\EventRow{78}{第4場面}{鳴き声}{ウサギ}{ウサギ鳴き声}{E3}{納得・理解を示し，移動開始（79/80）の前触れ}
	\EventRow{79}{第4場面}{足音}{ウサギ}{芝生の上を歩く}{4歩で一旦停止（タイミング設計）}{追従関係（距離を保ってついて行く）を時間構造で表現}
	\EventRow{80}{第4場面}{足音}{ウサギ}{草むらを歩く}{4歩で一旦停止（タイミング設計）}{追従関係（距離を保ってついて行く）を時間構造で表現}
	\EventRow{81}{第4場面}{鳴き声}{シカ}{baby-deer-calling-mama（息部）}{A2$\to$E3（息のみ提示）}{疑問でなく合図として鳴き声の単調さ回避を図る}
	\EventRow{82}{第4場面}{足音}{シカ（代替）}{馬が走る1}{C$\sharp$2（一定音量）}{ウサギに合わせて再出発し共同行動への移行を示す}
	\EventRow{83}{第4場面}{足音}{ウサギ}{芝生の上を歩く}{シカ停止後も約2秒継続+終盤ベロシティ上昇}{ウサギが遅れて到着し，川への到達を強調}
	\EventRow{84}{第4場面}{足音}{ウサギ}{草むらを歩く}{シカ停止後も約2秒継続+終盤ベロシティ上昇}{ウサギが遅れて到着し，川への到達を強調}
	\EventRow{85}{第4場面}{環境音}{環境}{風に揺れる草木1}{持続（背景）}{移動の場面を背景化し，自然環境の俯瞰感を作る}
	\EventRow{86}{第4場面}{環境音}{環境}{サンコクチョウのさえずり}{リバーブ付加}{移動シーンの単調さ回避／環境音としての広がりを付与}
	\EventRow{87}{第4場面}{環境音}{環境}{オオルリのさえずり1}{リバーブ付加}{移動シーンの単調さ回避／環境音としての広がりを付与}
	\EventRow{88}{第4場面}{環境音}{環境}{トンビの鳴き声}{リバーブ付加}{移動シーンの単調さ回避／環境音としての広がりを付与}
	\EventRow{89}{第4場面}{環境音}{環境}{川}{音量無音$\to$−8.0 dB（約14秒ランプ）}{川への接近を音量増加で表現}
	\EventRow{90}{第4場面}{鳴き声}{ウサギ}{ウサギ鳴き声}{G2$\to$（2秒）$\to$D3}{到達の納得・達成感を疑問と区別して提示し場面転換を示す}
	\EventRow{91}{第4場面}{鳴き声}{ハチドリ（主）}{（75の流用）ハチドリ鳴き声}{流用}{ウサギと対照的に疑問を残し，後続の態度差の導入}
	\EventRow{92}{第4場面}{鳴き声}{シカ}{baby-deer-calling-mama}{合図として配置}{自ら水を飲む行動と同様行動の促しを示す}
	\EventRow{93}{第4場面}{足音}{シカ（代替）}{馬が走る1}{短距離（2歩分）・一定音量}{水面までの近距離移動を示す}
	\EventRow{94}{第4場面}{効果音}{シカ（代替）}{dog-drinking-water-5}{C3$\to$A2}{シカの飲水を代替し，大型感を低ピッチで補う}
	\EventRow{95}{第4場面}{鳴き声}{シカ}{baby-deer-calling-mama}{E3+動作後1秒}{飲水終了の区切りと肯定的呼びかけを示す}
	\EventRow{96}{第4場面}{鳴き声}{ハチドリ（主）}{ハチドリ鳴き声（3連）}{連続発声+E3}{一時的な場の和らぎ／肯定的受容を示す（判断確定ではない）}
	\EventRow{97}{第4場面}{鳴き声}{ウサギ}{（90の流用）ウサギ鳴き声}{流用}{早い受容と行動移行の性格を補強}
	\EventRow{98}{第4場面}{鳴き声}{ウサギ}{ウサギ鳴き声}{C3（原音）}{感情でなく意思決定（次行動）を示し主体予告も担う}
	\EventRow{99}{第4場面}{足音}{ウサギ}{芝生の上を歩く}{短距離（3歩分）}{水面までの近距離移動を示す}
	\EventRow{100}{第4場面}{足音}{ウサギ}{草むらを歩く}{短距離（3歩分）}{水面までの近距離移動を示す}
	\EventRow{101}{第4場面}{鳴き声}{ウサギ}{ウサギ鳴き声}{G$\sharp$2}{飲水直前の準備完了・ためらいの少なさを示す}
	\EventRow{102}{第4場面}{効果音}{ウサギ}{dog-drinking}{D$\sharp$4}{小型動物として軽い飲水音を提示}
	\EventRow{103}{第4場面}{鳴き声}{ウサギ}{ウサギ鳴き声（2連）}{D3}{動作完了後の収束と満足・納得を表現}
	\EventRow{104}{第4場面}{鳴き声}{ウサギ}{（24基調）ウサギ鳴き声}{G2のみ（溜め）}{主人公を下ろす動作の前の溜め／力の小ささを示す}
	\EventRow{105}{第4場面}{鳴き声}{ハチドリ（主）}{ハチドリ鳴き声}{C3短音}{予測可能化した状況（驚きではない）を示す}
	\EventRow{106}{第4場面}{効果音}{ハチドリ（主）}{草むらを歩く}{A2（芝生レイヤなし）}{衝撃を抑えた下ろし方／柔らかい接地を表現}
	\EventRow{107}{第4場面}{鳴き声}{ハチドリ（主）}{ハチドリ鳴き声}{A2$\to$E3（2音目長め）}{意図未理解の疑問を示しつつ戸惑いは小さい}
	\EventRow{108}{第4場面}{鳴き声}{ウサギ}{ウサギ鳴き声（2連）}{E3}{前向きな声掛けで行動を促す}
	\EventRow{109}{第4場面}{鳴き声}{ハチドリ（主）}{ハチドリ鳴き声}{A3$\to$G3（短下降）}{上昇を使わず戸惑いを示す（高域内下降でニュアンス付け）}
	\EventRow{110}{第4場面}{足音}{ハチドリ（主）}{革靴で歩く}{E3}{物理再現より理解容易性を優先し，移動動作を明示}
	\EventRow{111}{第4場面}{鳴き声}{ハチドリ（主）}{ハチドリ鳴き声（短）}{短音}{思考の後に行動へ踏み出す判断成立を示す}
	\EventRow{112}{第4場面}{効果音}{ハチドリ（主）}{dog-drinking}{F$\sharp$4+間隔配置}{小型ゆえ高ピッチ化し，慎重に飲む（間を取る）態度を表現}
	\EventRow{113}{第4場面}{鳴き声}{ハチドリ（主）}{ハチドリ鳴き声}{D4$\to$C4$\to$B3$\to$A3（下降）}{飲水後の安堵感を段階下降で提示}
	\EventRow{114}{第4場面}{鳴き声}{ウサギ}{（108基調）ウサギ鳴き声}{E3$\to$D$\sharp$3（半音下げ）}{働きかけから受容・納得へ印象を微調整}
	\EventRow{115}{第4場面}{鳴き声}{シカ}{baby-deer-calling-mama}{B2}{周囲環境への反応として，低トーンで場面の気配を変える}
	\EventRow{116}{第4場面}{足音}{シカ（代替）}{馬が走る1}{1歩分のみ}{移動ではなくとっさの反応を示し，予兆として機能}
	\EventRow{117}{第4場面}{効果音}{環境}{風に揺れる草木2（流用）}{流用（強調）}{新たな存在（サル）出現の前触れを具体化}
	\EventRow{118}{第4場面}{鳴き声}{ハチドリ（主）}{ハチドリ鳴き声（上昇）}{F3}{これまでより強い動揺を上昇要素で表現}
	\EventRow{119}{第4場面}{鳴き声}{ウサギ}{ウサギ鳴き声}{極短・F3}{動揺より警戒として制御された反応を対照的に提示}
	\EventRow{120}{第4場面}{効果音}{サル}{草むらを歩く（流用）}{5つの異なるピッチを同時再生}{大型動物の飛び出しを音響密度で表現しダイナミクスを強調}
	\EventRow{121}{第4場面}{鳴き声}{サル}{monkey}{リバーブ付加}{距離・不確定な音像を作り，警戒対象として提示}
	\EventRow{122}{第4場面}{効果音}{サル}{芝生の上を歩く（流用）}{瞬間再生+リバーブ付加（約3秒遅れ）}{着地の衝撃と規模感を補完し存在感を強調}
	\EventRow{123}{第4場面}{効果音}{サル}{草むらを走る}{音量−17.4$\to$+3.5 dBランプ}{急接近を運動量の大きい素材と音量増加で表現}
	\EventRow{124}{第4場面}{鳴き声}{ハチドリ（主）}{ハチドリ鳴き声（2段）}{素材内で2回目高ピッチ}{警戒のためらい$\to$強い動揺への移行を示す}
	\EventRow{125}{第4場面}{鳴き声}{サル}{monkey（別素材）}{素材変更}{穏やかな発声で警戒対象とは異なる印象へ転換を作る}
	\EventRow{126}{第4場面}{効果音}{サル}{芝生の上を歩く（流用）}{瞬間再生（リバーブなし）}{走行停止の反動（遅れて接地）を示し近接を明示}
	\EventRow{127}{第4場面}{足音}{サル}{芝生の上を歩く}{音量−8.3$\to$0 dBランプ}{さらに接近する歩行を音量で表現}
	\EventRow{128}{第4場面}{鳴き声}{ハチドリ（主）}{ハチドリ鳴き声（弱$\to$強$\times$2）}{E3}{警戒継続を注意喚起的な音型で強調}
	\EventRow{129}{第4場面}{鳴き声}{ウサギ}{ウサギ鳴き声}{C2}{落ち着きを保った警戒を低ピッチで提示}
	\EventRow{130}{第4場面}{鳴き声}{サル}{monkey（121の素材）}{リバーブなし+D3}{近接かつ肯定的態度を示し印象転換を補強}
	\EventRow{131}{第4場面}{鳴き声}{ハチドリ（主）}{ハチドリ鳴き声（2連）}{G2+間を置く}{反射でなく思考的応答であることを時間構造で示す}
	\EventRow{132}{第4場面}{鳴き声}{サル}{monkey（軽快箇所）}{D3}{友好的態度を示し，印象の断絶を緩和して連続性を確保}
	\EventRow{133}{第4場面}{効果音}{サル}{衣擦れ1}{動作に同時付加}{草を渡す腕の動作を直感的に補完}
	\EventRow{134}{第4場面}{効果音}{サル}{collapsing-in-grass}{E3}{差し出された物が草であることと適度な質量感を提示}
	\EventRow{135}{第4場面}{鳴き声}{ハチドリ（主）}{（128と同箇所）ハチドリ鳴き声}{2回目開始で途中停止}{発声の未完結で判断停止＝困惑を聴覚的に表す}
	\EventRow{136}{第4場面}{鳴き声}{サル}{monkey（132と同箇所）}{E3}{意図伝達を強める積極的呼びかけを表現}
	\EventRow{137}{第4場面}{効果音}{サル}{衣擦れ1}{D$\sharp$3}{ジェスチャーとしての腕振りを活発・軽快に示す}
	\EventRow{138}{第4場面}{効果音}{サル}{collapsing-in-grass}{G$\sharp$3}{草の音を軽快化し，活発なジェスチャーの結果を補完}
	\EventRow{139}{第4場面}{鳴き声}{ハチドリ（主）}{ハチドリ鳴き声（短$\times$2）}{2連+B3}{慎重さを残した構え（行動前）を示す}
	\EventRow{140}{第4場面}{鳴き声}{ハチドリ（主）}{（139と同素材）ハチドリ鳴き声}{C4}{心理表現と区別し，動作に伴う掛け声として提示}
	\EventRow{141}{第4場面}{効果音}{ハチドリ（主）}{collapsing-in-grass（流用）}{D4を0.5秒$\to$F4短音}{保持$\to$引き離しの受け取り完了動作を音で表現}
	\EventRow{142}{第4場面}{効果音}{ハチドリ（主）}{リンゴをかじる（代替）}{間隔配置（3秒$\to$2秒）}{葉の咀嚼を代替し，慎重に緊張が解ける過程を示す}
	\EventRow{143}{第4場面}{効果音}{ハチドリ（主）}{飲む}{直後配置}{摂食の確定（飲み込み）として出来事を確定させる}
	\EventRow{144}{第4場面}{鳴き声}{ハチドリ（主）}{ハチドリ鳴き声（上昇）}{A3$\to$D4（上昇反復）}{行為完了後の安堵・満足を余韻として提示}
	\EventRow{145}{第4場面}{鳴き声}{サル}{monkey（132/136同系）}{C4}{意図達成への肯定を高ピッチで明確化}
	\EventRow{146}{第4場面}{鳴き声}{ウサギ}{ウサギ鳴き声}{B3}{サルに続く肯定反応としてまとまりを作る}
	\EventRow{147}{第4場面}{足音}{シカ（代替）}{馬が走る1}{1歩分のみ}{次の鳴き声の主体がシカであることを予告（識別補助）}
	\EventRow{148}{第4場面}{鳴き声}{シカ}{（95の流用）baby-deer-calling-mama}{流用+E3}{穏やかな肯定でキャラ差（抑制的反応）を表現}
	\EventRow{149}{第4場面}{鳴き声}{ハチドリ（主）}{ハチドリ鳴き声（弱$\to$強）}{A3短音+A3（2回目まで）}{弱い発声を基準化し強い発声を際立たせ，場面の締めとして元気さを提示}
\end{longtable}

\small
\clearpage
\twocolumn
\endinput
% =========================================
